\documentclass[11pt]{article}

\usepackage[T1]{fontenc}
\usepackage{amsmath,amssymb,amsthm,mathtools}
\usepackage[margin=1in]{geometry}
\usepackage[hidelinks]{hyperref}
\usepackage{microtype}

\newtheorem{theorem}{Theorem}[section]
\newtheorem{lemma}[theorem]{Lemma}
\newtheorem{proposition}[theorem]{Proposition}
\newtheorem{corollary}[theorem]{Corollary}
\theoremstyle{remark}
\newtheorem{remark}[theorem]{Remark}

\newcommand{\nuG}{\nu}
\newcommand{\aG}{a}
\newcommand{\F}{\mathbb F}
\newcommand{\Z}{\mathrm Z}
\newcommand{\CG}{\mathrm C}
\newcommand{\Span}{\operatorname{span}}
\newcommand{\Tr}{\operatorname{Tr}}

\title{Abelian Covers and Pairwise Noncommuting Sets in Groups:\\
       Structural Reductions, Self-Contained Bounds, and a Counterexample}
\author{Research draft}
\date{13 August 2026}

\begin{document}
\maketitle

\begin{abstract}
For a group $G$, let $\nuG(G)$ be the largest size of a pairwise
noncommuting set and let $\aG(G)$ be the least number of abelian subgroups
covering $G$.  Erd\H{o}s asked for the extremal function
$h(n)=\sup\{\aG(G):\nuG(G)\le n\}$.  We do not determine $h(n)$ and hence
do not claim a complete solution of Erd\H{o}s Problem~117.  We do give a
self-contained exact reduction from arbitrary groups to finite groups that
preserves the entire central-coset commutation graph; prove
$\nuG(\CG_G(x))\le\nuG(G)-2$ for every noncentral $x$, and consequently
$h(n)\le n h(n-2)$; determine $h(1)=h(2)=1$, $h(3)=3$, and $h(4)=4$; and
construct scalar symplectic groups $S(q,m)$ with
$\aG(S(q,m))=q^m+1$.  A complete calculation in $S(3,2)$ gives
$(\nuG,\aG)=(7,10)$, disproving the candidate formula considered here,
$h(n)=\max\{n,2^{\lfloor(n-1)/2\rfloor}+1\}$.  Exact computations provide
independent checks but are not dependencies of these proofs.
\end{abstract}

\section{Problem and status}

For a nonempty group $G$, define
\[
 \nuG(G)=\max\{|X|:xy\ne yx\text{ for all distinct }x,y\in X\}
\]
when this maximum is finite, and let $\aG(G)$ be the minimum cardinality of
a family of abelian subgroups whose union is $G$.  The object of the problem
is
\[
 h(n)=\sup\{\aG(G):G\text{ a group and }\nuG(G)\le n\}.
\]
Erd\H{o}s's Problem~26 asks for this function and records two-sided
exponential estimates attributed to earlier work
\cite[p.~8]{erdos_1997}. Neumann's Theorem~6 proves the underlying
arbitrary-group central-by-finite theorem \cite[p.~470]{neumann_1976}, and
Faber, Laver and McKenzie give the graph-cover correspondence and qualitative
equivalences \cite[p.~933 and Theorem~3, pp.~936--937]
{faber_laver_mckenzie_1978}. Pyber's publisher abstract records, for finite
groups, an absolute constant $c$ such that
$[G:\Z(G)]\le c^{\nuG(G)}$ \cite{pyber_1987}; its full proof was not audited.
We use none of those results as a
load-bearing step below: the qualitative reduction and all displayed
quantitative bounds are proved here.

The exact asymptotic growth of $h$ remains unresolved in this manuscript.
In particular, we prove only
\[
 \liminf_{n\to\infty}h(n)^{1/n}\ge\sqrt2
\]
and a superexponential self-contained upper bound.  The fixed-base upper
bound reported by Pyber is historically stronger, but its proof is not
reproduced here.

\section{The noncommuting graph and central-coset compression}

Let $\Gamma_G$ be the graph with vertex set $G$, where distinct vertices are
adjacent when they do not commute.

\begin{proposition}[Graph/group dictionary]\label{prop:dictionary}
For every group $G$,
\[
 \omega(\Gamma_G)=\nuG(G),\qquad \chi(\Gamma_G)=\aG(G),
\]
with the natural cardinal interpretation before finiteness is known.
\end{proposition}

\begin{proof}
The first equality is the definition.  Every abelian-subgroup cover gives a
proper coloring by assigning each element one subgroup containing it.
Conversely, each color class is a pairwise commuting set and therefore
generates an abelian subgroup.  These generated subgroups form a cover.
\end{proof}

Define $\Delta_G$ on $G/\Z(G)$, retaining the identity coset, by joining
$x\Z(G)$ and $y\Z(G)$ exactly when $xy\ne yx$.

\begin{proposition}[Exact central-coset compression]\label{prop:compression}
The graph $\Delta_G$ is well defined and
\[
 \omega(\Delta_G)=\nuG(G),\qquad\chi(\Delta_G)=\aG(G).
\]
\end{proposition}

\begin{proof}
For $z,w\in\Z(G)$, the elements $xz$ and $yw$ commute exactly when $x$ and
$y$ commute.  Thus $\Gamma_G$ is obtained from $\Delta_G$ by replacing every
vertex with an independent set of cardinality $|\Z(G)|$, with complete or
empty bipartite connections between fibers according to the base graph.
A clique uses at most one point in each fiber and every clique of the base
lifts.  A coloring of the base pulls back, while restricting a coloring to
a transversal gives the reverse chromatic inequality.
\end{proof}

The graph $\Delta_G$ is generally \emph{not} the noncommuting graph of
$G/\Z(G)$: quotient commutation asks only whether $[x,y]\in\Z(G)$.  For
example, $D_8/\Z(D_8)$ is abelian while $\Delta_{D_8}$ consists of an
isolated vertex and a triangle.

\section{A self-contained finite model}

We first obtain finite center index without invoking the qualitative
literature.  Write $R(s,t)$ for the Ramsey number.

\begin{lemma}[Uniform centralizer index]\label{lem:bfc}
If $\nuG(G)=n<\infty$, then
\[
 [G:\CG_G(g)]\le R(n+1,n+1)-1
\]
for every $g\in G$.
\end{lemma}

\begin{proof}
If the conjugacy class of $g$ had at least $R(n+1,n+1)$ elements, choose a
set $T$ of that size giving distinct conjugates $t^{-1}gt$.  Ramsey's theorem
cannot return $n+1$ pairwise noncommuting members of $T$, so it returns
commuting $t_1,\dots,t_{n+1}$.  For distinct commuting $u,v\in T$, equality
$(gu)(gv)=(gv)(gu)$ is equivalent to $u^{-1}gu=v^{-1}gv$.  The latter is
false, so $gT$ contains $n+1$ pairwise noncommuting elements, a contradiction.
\end{proof}

\begin{lemma}[A maximum clique detects the center]\label{lem:center}
If $X=\{x_1,\dots,x_n\}$ is a maximum pairwise noncommuting set, then
\[
 \bigcap_{i=1}^n\CG_G(x_i)=\Z(G).
\]
\end{lemma}

\begin{proof}
Only the reverse inclusion needs proof.  Suppose $c$ centralizes $X$ but is
not central, and choose $y$ not commuting with $c$.  Replace $x_i$ by
$x_ic$ precisely when $x_i$ commutes with $y$.  Multiplication by the common
element $c$, which centralizes every $x_i$, preserves all pairwise
noncommutation inside $X$, while each modified element fails to commute with
$y$.  Adjoining $y$ gives a larger noncommuting set.
\end{proof}

Lemmas~\ref{lem:bfc} and \ref{lem:center} show that $\Z(G)$ is a finite
intersection of finite-index subgroups, hence has finite index.

\begin{theorem}[Exact finite commutation model]\label{thm:finite-model}
If $[G:\Z(G)]<\infty$, there is a finite group $K$ for which
\[
 \Delta_K\cong\Delta_G,\qquad \nuG(K)=\nuG(G),\qquad \aG(K)=\aG(G).
\]
Consequently, for every $n\ge1$, the supremum defining $h(n)$ is attained
by a finite group.
\end{theorem}

\begin{proof}
Let $Z=\Z(G)$, choose a finite transversal for the central subgroup
$T=\{t_1=1,t_2,\dots,t_q\}$, put $H=\langle T\rangle$ and $C=H\cap Z$.
Then $H/C\cong G/Z$ and $C$ has finite index in the finitely generated group
$H$.  Schreier rewriting shows that $C$ is finitely generated; it is
abelian because $C\le Z$.  Let
\[
 D=\{[t_i,t_j]:[t_i,t_j]\ne1\text{ and }[t_i,t_j]\in C\}.
\]
This is finite.  A finitely generated abelian group is residually finite on
finite sets, so there is a finite-index $N\le C$ with $N\cap D=\varnothing$.
For completeness, after writing $C\cong\mathbb Z^r\oplus F$, one may take
$N=M\mathbb Z^r\oplus0$, where $M$ is divisible by the exponent of $F$ and
larger than the absolute value of every nonzero free coordinate appearing
in $D$.

Set $K=H/N$.  It is finite.  Since $C$ is central, the commutator of
$t_ic$ and $t_jd$ is $[t_i,t_j]$.  A nontrivial such commutator is not in
$N$: either it lies outside $C$, or it belongs to $D$.  Thus the quotient
creates no new commuting pair between transversal cosets.  Moreover
$Z(K)=C/N$: if $t_icN$ is central, then every $[t_i,t_j]$ lies in $N$;
commutation preservation forces all of them to be trivial. Thus $t_i$
centralizes $H$. Conversely this condition plainly makes the coset central.
Since $G=HZ$, it is equivalent to $t_i\in Z$, which forces $i=1$.
Therefore $K/Z(K)\cong H/C\cong G/Z$, preserving every edge of $\Delta$.
Proposition~\ref{prop:compression} proves the invariant equalities.

For the last assertion, first apply Lemmas~\ref{lem:bfc} and
\ref{lem:center}.  The bounds below show that $h(n)$ is a finite integer.
A bounded nonempty set of integers attains its supremum, and the preceding
construction replaces any extremizer by a finite one.
\end{proof}

\section{Universal upper bounds}

\begin{lemma}[Two-step centralizer drop]\label{lem:drop}
If $x\notin\Z(G)$ and $\nuG(G)<\infty$, then
\[
 \nuG(\CG_G(x))\le\nuG(G)-2.
\]
\end{lemma}

\begin{proof}
Let $a_1,\dots,a_\ell$ be pairwise noncommuting elements of $\CG_G(x)$,
and choose $y$ with $xy\ne yx$.  Put
\[
 b_i=\begin{cases}a_i,&a_iy\ne ya_i,\\ xa_i,&a_iy=ya_i.\end{cases}
\]
Because every $a_i$ commutes with $x$, multiplying selected $a_i$'s by $x$
preserves pairwise noncommutation among them.  If $a_i$ does not commute with
$y$, it cannot commute with $xy$.  If it does commute with $y$, then $xa_i$
commutes with neither $y$ nor $xy$, since either equality would imply
$xy=yx$.  Finally $y$ and $xy$ do not commute.  Hence
$\{b_1,\dots,b_\ell,y,xy\}$ is pairwise noncommuting, proving
$\ell+2\le\nuG(G)$.
\end{proof}

\begin{theorem}[Recursive upper bound]\label{thm:recurrence}
For $n\ge3$,
\[
 h(n)\le n h(n-2).
\]
In particular, for $m\ge1$,
\[
 h(2m+1)\le(2m+1)!!,\qquad
 h(2m)\le\frac{(2m)!!}{2}=2^{m-1}m!.
\]
\end{theorem}

\begin{proof}
If $G$ is abelian, then $\aG(G)=1$ and there is nothing to prove. Otherwise,
if $S$ is a maximum noncommuting set, then maximality gives
$G=\bigcup_{x\in S}\CG_G(x)$.  Each $x\in S$ is noncentral.  If
$|S|=r\le n$, Lemma~\ref{lem:drop} and monotonicity of $h$ give
\[
 \aG(G)\le\sum_{x\in S}\aG(\CG_G(x))
 \le r h(r-2)\le n h(n-2).
\]
A nonabelian group contains the pairwise noncommuting triple $x,y,xy$;
hence $h(1)=h(2)=1$.  Iteration proves the stated bounds.
\end{proof}

We also record a second, independent quantitative route.  It is weaker than
Theorem~\ref{thm:recurrence} as a bound for $h$, but gives useful structural
control of conjugacy classes.

\begin{proposition}[Conjugacy and center index]\label{prop:conjugacy}
If $\nuG(G)=n<\infty$, then every conjugacy class has size at most $4n^2$ and
\[
 [G:\Z(G)]\le(4n^2)^n,\qquad \aG(G)\le(4n^2)^n-1.
\]
\end{proposition}

\begin{proof}
First suppose $G$ finite.  List its conjugacy classes $D_1,\dots,D_s$ in
nondecreasing order, and choose minimal $r$ with
$|Y|:=\sum_{i\le r}|D_i|>|G|/2$.  The complement of
$D_1\cup\cdots\cup D_{r-1}$ has size at least $|G|/2$. Choose an
inclusion-maximal clique in the induced graph; it has at most $n$ elements,
and its centralizers cover that complement. Some such
centralizer has size at least $|G|/(2n)$, whence $|D_r|\le2n$.
For each $g$, the sets $Y$ and $gY$ intersect, so $g\in D_iD_j^{-1}$ for
some $i,j\le r$.  This product is conjugacy invariant; consequently
$|g^G|\le|D_i||D_j|\le4n^2$.

For arbitrary $G$, Theorem~\ref{thm:finite-model} produces a finite $K$
with the same central-coset graph.  The closed non-neighborhood of a coset
records its centralizer cosets, so corresponding elements have equal
centralizer index.  The finite bound transfers.  Finally apply
Lemma~\ref{lem:center} and the elementary inequality for an intersection of
finite-index subgroups:
\[
 [G:\Z(G)]\le\prod_{i=1}^n[G:\CG_G(x_i)]\le(4n^2)^n.
\]
If $G$ is nonabelian, one abelian subgroup
$\langle t,\Z(G)\rangle$ for each nonidentity central coset covers $G$,
giving the last inequality. If $G$ is abelian, the inequality follows
directly from $\aG(G)=1$.
\end{proof}

\begin{corollary}[Small exact values]\label{cor:small}
\[
 h(1)=h(2)=1,\qquad h(3)=3,\qquad h(4)=4.
\]
\end{corollary}

\begin{proof}
The first two values were proved above.  Theorem~\ref{thm:recurrence} gives
$h(3)\le3$ and $h(4)\le4$.  The group $D_8$ has central-coset graph an
isolated vertex plus a triangle, so $(\nuG,\aG)=(3,3)$.  In $S_3$, the
three transpositions and either nonidentity $3$-cycle form a clique of size
four, while the three reflection subgroups and the subgroup of order three
form a cover of size four.
\end{proof}

\section{Scalar symplectic groups}

Let $q$ be a prime power and $m\ge1$.  On
\[
 S(q,m)=\F_q^m\times\F_q^m\times\F_q
\]
define
\[
 (x,y,z)(x',y',z')=(x+x',y+y',z+z'+x\mathbin\cdot y').
\]
Bilinearity gives associativity and
$(x,y,z)^{-1}=(-x,-y,-z+x\cdot y)$.  The center is the last coordinate,
and commutation on $V=\F_q^m\oplus\F_q^m$ is orthogonality for
\[
 B((x,y),(x',y'))=x\cdot y'-x'\cdot y.
\]
Let $\pi(q,m)$ be the maximum number of projective points of $V$ that are
pairwise nonorthogonal.

\begin{theorem}[Symplectic parameters]\label{thm:symplectic}
For every prime power $q$ and $m\ge1$,
\[
 \nuG(S(q,m))=\pi(q,m),\qquad \aG(S(q,m))=q^m+1.
\]
Moreover $\pi(2,m)=2m+1$.
\end{theorem}

\begin{proof}
The first equality follows because nonzero scalar multiples are orthogonal
and have identical orthogonality relations to every other vector.

The image of an abelian subgroup in $V$ is an additive set on which $B$
vanishes; its $\F_q$-span is totally isotropic.  Such a subspace has
dimension at most $m$, so each cover member contains at most $q^m-1$ of the
$q^{2m}-1$ nonzero central cosets.  Hence $\aG\ge q^m+1$.  For the reverse
inequality, identify $V$ symplectically with $\F_{q^m}^2$ endowed with
\[
 ((u,v),(u',v'))\longmapsto\Tr_{\F_{q^m}/\F_q}(uv'-u'v).
\]
The $q^m+1$ one-dimensional $\F_{q^m}$-subspaces are totally isotropic,
meet pairwise only in zero, and partition the nonzero vectors.  Their full
preimages give an abelian cover.

To justify the identification explicitly, the trace map of a finite-field
extension is a nonzero $\F_q$-linear functional. If $u\ne0$, choose $w$ with
$\Tr(w)\ne0$ and take $v'=u^{-1}w$. Hence the trace pairing, and therefore
the displayed alternating form, is nondegenerate. The symplectic-basis
induction identifies any two nondegenerate alternating forms of dimension
$2m$.

When $q=2$, a pairwise nonorthogonal set $v_1,\dots,v_r$ has Gram matrix
with zero diagonal and one off the diagonal.  Its rank over $\F_2$ is $r$
for even $r$ and $r-1$ for odd $r$, so $r\le2m+1$.  Conversely the
$2m\times2m$ such matrix is nondegenerate alternating and hence isometric
to $B$. Under such an isometry, the standard ordered basis gives vectors
$v_1,\dots,v_{2m}$ with that Gram matrix; put
$v_{2m+1}=\sum_i v_i$; all pairings are one.  Thus $\pi(2,m)=2m+1$.
\end{proof}

\begin{corollary}[Exponential lower bound]\label{cor:lower}
For $m\ge1$,
\[
 h(2m+1)\ge2^m+1,\qquad
 \liminf_{n\to\infty}h(n)^{1/n}\ge\sqrt2.
\]
\end{corollary}

\begin{proof}
The first inequality follows from Theorem~\ref{thm:symplectic} with $q=2$.
Monotonicity transfers it from $n=2m+1$ to the following even index, and
taking roots along the two interlaced subsequences gives the limit inferior.
\end{proof}

\section{A counterexample to the binary exact candidate}

The binary family suggests the tempting formula
\[
 h(n)=\max\{n,2^{\lfloor(n-1)/2\rfloor}+1\}.\tag{$*$}
\]
It already gives the wrong value $2$ at $n=1,2$.  More substantially, its
restriction to $n\ge3$ is also false.

\begin{theorem}\label{thm:counterexample}
The group $S(3,2)$ has order $3^5$ and
\[
 (\nuG(S(3,2)),\aG(S(3,2)))=(7,10).
\]
Consequently $h(7)\ge10$ and $h(8)\ge10$, whereas $(*)$ gives $9$ at both
indices.
\end{theorem}

\begin{proof}
The cover number is Theorem~\ref{thm:symplectic}.  We prove the clique
number in the interleaved coordinates $v=(a,b,c,d)$ with
$B(v,v')=ab'-ba'+cd'-dc'$.  The seven vectors
\[
\begin{gathered}
(1,0,0,0),(0,1,0,0),(1,1,0,0),\\
(-1,1,0,1),(-1,1,1,0),(-1,1,1,1),(-1,1,1,-1)
\end{gathered}
\]
are pairwise nonorthogonal: the last four use the four projective directions
of $\F_3^2$ in their final coordinates, and their identical first
coordinates contribute zero to mutual pairings; pairings with the first
three are immediate.

For the upper bound, there is nothing to show for fewer than two points.
Otherwise rescale and order two nonorthogonal vectors so their pairing is
one, then extend them to a symplectic basis, normalizing them to
$(1,0,0,0)$ and $(0,1,0,0)$.
Every remaining projective point has a unique representative
$(\epsilon,1,u)$ with $\epsilon\in\{1,-1\}$ and $u\in\F_3^2$.  Let $P$ and
$M$ be the $u$'s belonging respectively to $\epsilon=1$ and $-1$.  They
satisfy
\[
 \det(p,p')\ne0,\quad\det(m,m')\ne0,\quad\det(p,m)\ne1.\tag{1}
\]
Each of $P,M$ uses at most one vector in each of the four directions.  If
one has size at most one, then immediately $|P|+|M|\le5$. If both have size
at least two, reorder two members of $P$ so their determinant is $-1$ and
apply an element of $\mathrm{SL}_2(3)$ sending them to
$(0,1),(1,0)$. Condition (1) restricts $M$ to
\[
 (0,-1),\quad(1,0),\quad(1,-1).
\]
If $|P|=2$, this gives $|P|+|M|\le5$. A third member $p=(x,y)$ of $P$ is
one of the following four vectors; the final column lists candidates
excluded by $\det(p,m)\ne1$:
\[
\begin{array}{c|c}
p&\text{excluded candidates}\\ \hline
(1,1)&(1,-1)\\
(1,-1)&(1,0)\\
(-1,1)&(0,-1)\\
(-1,-1)&(0,-1),(1,0)
\end{array}
\]
Thus $|M|\le2$. If $|P|=4$, its two additional directions have
representatives in $\{(1,1),(-1,-1)\}$ and
$\{(1,-1),(-1,1)\}$; substituting the four choices leaves at most one
candidate in $M$.  Thus always $|P|+|M|\le5$, and the original set has at
most $2+5=7$ points.
\end{proof}

\section{Direct products and why they do not settle the problem}

For graphs $X,Y$, their OR product $X\vee Y$ joins two product vertices
when at least one coordinate is adjacent.

\begin{proposition}\label{prop:products}
For groups $G,H$,
\[
 \Gamma_{G\times H}=\Gamma_G\vee\Gamma_H.
\]
For every nonempty finite graph $X$,
\[
 \lim_{k\to\infty}\chi(X^{\vee k})^{1/k}=\chi_f(X),\qquad
 \lim_{k\to\infty}\omega(X^{\vee k})^{1/k}=\Theta(\overline X),
\]
where $\chi_f$ is fractional chromatic number and $\Theta$ is Shannon
capacity.
\end{proposition}

\begin{proof}
Two product-group elements fail to commute exactly when one coordinate pair
fails to commute.  Fractional chromatic number is multiplicative under OR
product by the primal and dual set-cover linear programs.  The standard
greedy set-cover inequality
$\chi(Y)\le(1+\log|V(Y)|)\chi_f(Y)$ removes only a subexponential factor
from powers.  Finally
$\overline{X\vee Y}=\overline X\boxtimes\overline Y$, so the clique sequence
is exactly the independence sequence defining $\Theta(\overline X)$.
\end{proof}

Thus powers of any fixed finite nonabelian group have both $\aG(G^k)$ and
$\nuG(G^k)$ exponential in $k$, making $\aG(G^k)$ only a power of
$\nuG(G^k)$.  Fixed-seed direct powers cannot produce a lower bound
$h(n)\ge c^n$; the varying scalar commutator geometry above can.

\section{Computational corroboration}

The repository contains dependency-free group implementations, exact
branch-and-bound clique search, exact DSATUR coloring, generated abelian
subgroup covers, an independent maximal-abelian set-cover formulation, and
certificate verification.  The structural theorem above for $S(3,2)$ was
discovered computationally and then proved independently.  The computation
excludes an eighth projective clique and colorings with at most nine colors,
and verifies a ten-subgroup symplectic-spread cover.  Full saved-data scope,
hashes, and commands are recorded in Appendix~\ref{app:computation}.

\section{What remains open}

The manuscript does not determine $h(n)$ beyond $n=4$, does not determine
the optimal exponential base, and does not prove that such a base exists.
The most concrete geometric subproblem is to control $\pi(q,m)$, the maximum
partial-ovoid size in a symplectic polar space, relative to the exact cover
number $q^m+1$.  More generally one must decide whether odd-characteristic
scalar forms, vector-valued commutator maps, solvable groups, or semisimple
groups can improve the efficiency achieved by the constructions above. An
upper bound near $2^{n/2}$ would require substantial new structure beyond
the recurrence in Theorem~\ref{thm:recurrence}.

Erd\H{o}s--Straus report that Isaacs obtained a factorial upper estimate and
finite examples with exponential cover size
\cite[p.~311]{erdos_straus_1976}. Later sources report the extraspecial
$2$-group construction and trace it through Bertram to an Isaacs personal
communication \cite[Example~4.5, p.~488]
{abdollahi_akbari_maimani_2006}\cite[ref.~6]{bertram_1983}. These are reports
rather than a checked primary proof. I\c{s}{\i}k, Theorem~12, pp.~6--7,
proves $\nuG(S(2,m))=2m+1$, but not the cover lower bound
\cite{isik_2005}. All claims in this paper that are needed for its
mathematical conclusions have instead been proved above.

% Reproducible computational methods and certificates will be documented here.


\bibliographystyle{plain}
\bibliography{main}

\end{document}
